\documentclass[acmsmall,screen,review,anonymous]{acmart}

\settopmatter{printfolios=true,printccs=false,printacmref=false}
\renewcommand\footnotetextcopyrightpermission[1]{}
\pagestyle{plain}

\usepackage{mathtools}
\usepackage{amsmath}
\usepackage{amsfonts}
\usepackage{amsthm}
\usepackage{booktabs}
\usepackage{enumitem}
\usepackage{xspace}

\ccsdesc[500]{Theory of computation~Program verification}
\ccsdesc[500]{Theory of computation~Equational logic and rewriting}
\ccsdesc[300]{Hardware~Quantum computation}

\keywords{Agda, quantum circuits, normal forms, presentations, Clifford group, symplectic group}

\newcommand{\Agda}{Agda\xspace}
\newcommand{\Zp}{\mathbb{Z}_p}
\newcommand{\Sp}{\mathrm{Sp}}
\newcommand{\Pauli}{\mathcal{P}}
\newcommand{\Cliff}{\mathcal{C}}
\newcommand{\NF}{\mathsf{NF}}
\newcommand{\Word}{\mathsf{Word}}
\newcommand{\Rel}{\mathsf{Rel}}
\newcommand{\code}[1]{\texttt{#1}}
\newcommand{\codepath}[1]{\nolinkurl{#1}}
\newcommand{\todo}[1]{\textbf{TODO:} #1}

\title{Certified Normal Forms and Complete Relations for Circuit Presentations in Agda}

\author{Anonymous Author(s)}
\affiliation{\institution{Anonymous Institution}}

\begin{abstract}
Completeness proofs for circuit rewrite systems often rely on long
calculations with generators, relations, normal forms, and case analyses over
families of circuits. These calculations are exactly the sort of argument that
is easy to get wrong on paper: a single missing commutation case or a mistaken
field identity can invalidate the claimed presentation. We describe an \Agda
framework for formalizing circuit presentations as words modulo congruence
relations, proving normal-form theorems for those presentations, and deriving
complete rewrite systems by verified normalization.

The central case study is a parameterized formalization of the symplectic group
\(\Sp(2n,\Zp)\), where \(n\) is arbitrary and \(p\) is an odd prime. We define
Pauli vectors, the symplectic form, the action of the generators \(S\), \(H\),
and \(CZ\), an inductive family of normal forms, implementability of every
symplectic linear map by a normal form, and injectivity of normal forms. We then
formalize the proof pattern that reduces arbitrary words to normal form by
pushing generators through normal boxes. This yields a certified route from
relations to completeness: two words are equal precisely when their normal forms
coincide.

Beyond the symplectic development, the framework supports semidirect products,
group extensions, coset normal forms, and amalgamated products. We use these
constructions to derive Clifford-mod-scalar relations from symplectic and Pauli
relations, simplify the Clifford presentation, explain how scalar phases can be
added as a general group extension, and verify several independent
amalgamation examples. The result is both a mechanized proof of a substantial
qudit Clifford presentation and a reusable library for certified algebraic
reasoning about circuit calculi.
\end{abstract}

\begin{document}
\maketitle

\section{Introduction}

Equational theories for circuits are useful only when their rule sets are
complete: if two circuits denote the same operation, the equations should be
able to transform one circuit into the other. Completeness is usually proved by
choosing a normal form, showing that every circuit reduces to one, and proving
that equal denotations have equal normal forms. This proof strategy is
conceptually clean, but in quantum and reversible circuit calculi the actual
proof is computationally dense. Even for well-structured fragments, the proof
requires many parameterized relations, many interactions between gates on
different wires, and many identities in finite rings or fields.

This paper describes a proof-assistant approach to these completeness proofs.
Instead of performing the relation reductions by hand, we formalize
presentations, words, normal forms, and proof-generating normalization in
\Agda~\cite{Agda}. The immediate payoff is reliability: the proof assistant
checks every algebraic step used to derive a box relation or simplify a
presentation. The broader payoff is reuse: once presentations and normal forms
are expressed abstractly, the same proof patterns apply to semidirect products,
coset normal forms, group extensions, and amalgamated products.

The main mathematical development is a mechanization of normal forms and
complete relations for \(\Sp(2n,\Zp)\), uniformly in the number of qudits \(n\)
and the odd prime dimension \(p\). The generators are the symplectic images of
the qudit Clifford gates \(S\), \(H\), and \(CZ\). The normal forms are built
from inductively defined boxes. Each box has a semantic action on Pauli vectors,
and each step of the normal-form proof is checked against the symplectic form.
The formalization proves that every symplectic linear transformation preserving
the form is implemented by a normal form, and that normal forms are unique.

The second focus is the path from the symplectic presentation to Clifford
presentations. The Clifford group modulo scalars decomposes into a Pauli layer
and a symplectic layer. We express this as a semidirect product presentation:
the Pauli presentation forms the normal subgroup, the symplectic presentation
acts by conjugation, and the combined relations specify how symplectic
generators move past Pauli generators. We then simplify the resulting
Clifford-mod-scalar relations by pushing Pauli components to canonical
positions and proving an identity isomorphism between the original and
simplified presentations. Finally, scalar phases are added using a more general
group-extension construction: quotient relators may lift with scalar
corrections, unlike the split semidirect-product case.

\paragraph{Contributions.}
The paper makes the following contributions.

\begin{itemize}[leftmargin=*]
\item We give a reusable \Agda representation of circuit presentations as word
  setoids: binary-tree words over generators, congruence closures of relations,
  and checked equational reasoning.
\item We formalize generic normal-form interfaces and construction principles,
  including direct products, semidirect products, coset normal forms,
  amalgamation, and group extensions.
\item We mechanize a normal-form theorem for \(\Sp(2n,\Zp)\), parameterized by
  \(n\) and an odd prime \(p\), including implementability and uniqueness of
  normal forms.
\item We formalize the completeness proof pattern ``push generators through
  boxes, then merge into normal form'', isolating the reusable normalization
  induction from the box-specific calculations.
\item We derive Clifford-mod-scalar relations from Pauli and symplectic
  relations by semidirect product, prove a simplified presentation equivalent,
  and explain how to add scalar phases by a general group-extension recipe.
\item We demonstrate the generality of the framework on amalgamation examples:
  single-qubit Clifford+\(T\), single-qutrit Clifford+\(T\), and the group of
  \(3\times3\) unitary matrices over \(\mathbb{Z}[1/2,i]\).
\end{itemize}

\section{Circuit Presentations in Agda}
\label{sec:framework}

The framework begins with a deliberately small syntactic representation. For a
set \(X\) of generators, circuits are words:
\[
  w ::= [x] \mid \epsilon \mid w \cdot w.
\]
The implementation uses binary trees rather than lists. This is not merely a
programming convenience: in \Agda, list concatenation forgets the exact operands
used to build the list, while a binary tree preserves the local shape needed by
congruence proofs. Associativity and unit laws are then added as relations,
recovering the usual monoid behavior at the setoid level.

A presentation is a relation \(\Gamma : \Word(X) \to \Word(X) \to \mathsf{Set}\).
The framework forms its congruence closure \(w \approx_\Gamma v\), generated by
reflexivity, symmetry, transitivity, congruence for composition, associativity,
unit laws, and the axioms in \(\Gamma\). This definition appears in
\code{Presentation/Base.agda}. Standard \Agda setoid reasoning is then used to
write checked derivations that look like ordinary equational proofs.

\subsection{Normal-Form Interfaces}

The main abstraction is \code{NFProperty}. A presentation has an
\code{NFProperty} when there is a type \(\NF\), a normalization function
\(\mathsf{nf} : \Word(X) \to \NF\), a proof that equivalent words have equal
normal forms, and a proof that equal normal forms imply word equivalence. The
stronger \code{NFProperty'} additionally stores an inverse map from normal-form
data back to words. These interfaces turn equality of circuits into equality of
normal-form data, which is often decidable or structurally simpler.

\todo{Insert a small excerpt of \code{NFProperty} and \code{NFProperty'} from
\codepath{Presentation/Properties.agda}, simplified for exposition.}

\subsection{Combinators for Presentations}

The library contains construction principles that transport normal-form
properties across larger presentations. Direct products combine independent
presentations. Semidirect products add cross-relations expressing an action of
one presentation on another. Coset normal forms extend a subgroup presentation
by a finite or structured set of representatives. Amalgamation combines two
presentations sharing a common subpresentation. Group extensions add a normal
subgroup presentation, a quotient presentation, corrected lifts of quotient
relations, and conjugation data.

These constructions are not merely organizational. They are the reason the
Clifford and example developments do not need to reprove normal-form machinery
from scratch. Each case supplies the local data required by the generic theorem,
and \Agda checks that the data respects the defining relations.

\section{The Symplectic Case Study}
\label{sec:symplectic}

Let \(p\) be an odd prime. The formalization represents a Pauli vector on \(n\)
qudits as a vector of pairs in \(\Zp \times \Zp\). For a single qudit, the
symplectic form is
\[
  \langle (a,b),(c,d)\rangle = -a d + c b,
\]
and the \(n\)-qudit form is the sum of these contributions across the vector.
The module \code{N.Pauli} defines the Pauli vector type, the form, and basic
lemmas such as antisymmetry and the behavior of the identity Pauli vector.

The generator set is indexed by arity. The single-qudit generators \(S\) and
\(H\) live at every positive arity; the two-qudit generator \(CZ\) lives at
arity at least two; and a lift operation shifts a generator to a higher wire.
Derived words such as \(CX\), \(SWAP\), and multipliers are defined from these
basic generators. The presentation \code{N.Symplectic} contains the
parameterized relations used for the symplectic group.

\subsection{Normal Boxes}

The normal forms are built inductively from boxes. Informally, an \(n\)-qudit
normal form consists of an \((n-1)\)-qudit normal form lifted to the upper wires
and an \(LM\) layer acting on the newly exposed wire. The boxes \(A,B,D,E,L,M\)
encode structured fragments that eliminate or normalize the first nontrivial
part of a symplectic tableau. In \Agda these boxes are data types indexed by
arity, with an interpretation as words.

\todo{Add a figure for the inductive normal form, adapted from Figure 3 of
\codepath{doc/qudit-quantum-June12.pdf} if permissions/attribution are
settled, or redraw as a simple schematic.}

\subsection{Implementability and Uniqueness}

The semantic action of a word on Pauli vectors is defined generator by
generator and extended compositionally. The formalization proves that this
action is linear and preserves the symplectic form. Conversely, for every
linear function preserving the symplectic form, \code{N.NF} constructs a normal
form whose action is inverse to that function. Applying this construction to
the inverse of a word yields a normal form implementing the word itself.

Uniqueness is proved by showing that distinct normal forms have distinct Pauli
actions. The proof is inductive. At the outer layer, equality of actions implies
equality of the head behavior of the \(LM\) component; a head-injectivity lemma
identifies the \(LM\) boxes. Surjectivity of the tail action then reduces the
problem to the inner normal forms. This is the core of
\code{N.NF-Inj.agda}.

\begin{theorem}[Normal forms for \(\Sp(2n,\Zp)\), mechanized]
For every \(n\) and every odd prime \(p\), every word over \(S\), \(H\), and
\(CZ\) has a normal form with the same Pauli action. Moreover, two normal forms
with the same action are equal.
\end{theorem}

\todo{Replace this theorem statement with the exact names and dependencies once
the trusted artifact slice is finalized.}

\section{Completeness by Pushing Through Boxes}
\label{sec:completeness}

The completeness proof separates box-specific algebra from a generic
normalization induction. The box-specific part proves that each generator can be
pushed through each normal box, producing a new box and some residual word. For
example, pushing \(S\) through an \(E\)-box updates the parameter:
\[
  E_b \cdot S \approx E_{b-1}.
\]
The theorem \codepath{BoxRelations.One.E<-S}, backed by
\codepath{N.BR.One.E.lemma-single-qupit-br-E}, is the smallest example of a
certified box reduction: it uses the single-qudit \(S\)-power lemma and
arithmetic in \(\Zp\) to prove
\[
  [b]^e \cdot S \approx [b-1]^e.
\]
At the other end of the scale, \codepath{BoxRelations.Three.DD<-CZ}, backed by
\codepath{N.BR.Three.DD-CZ.lemma-dir-and-vd'}, pushes a \(CZ\) through a two-entry
\(DD\) box:
\[
  [vd]^{vd} \cdot CZ \approx dir^\uparrow \cdot [vd']^{vd}.
\]
The nonzero/nonzero cases require substantial reassociation, commutation, and
movement of three-qudit gadgets. These examples illustrate why the proof
assistant matters: the statement of the box relation is compact, but the checked
derivation can be long and easy to miscopy by hand.

The generic part is implemented in \code{Symplectic-Completeness.agda}. The
interface is a single operation:
\[
  \mathsf{cosetUpdate}(lm,g)
    = (lm',w,\; [lm]\cdot[g] \approx w^\uparrow\cdot[lm']).
\]
Given this operation, the rest of completeness is structural. To merge a gate
into a normal form \((nf,lm)\), push the gate through \(lm\), then recursively
merge the residual word \(w\) into the inner normal form \(nf\). To normalize a
word, repeat this merge operation over the word structure.

\begin{theorem}[Completeness by normalization, mechanized skeleton]
Assuming the box-relation chaining operation \(\mathsf{cosetUpdate}\), every
word appended to a normal form rewrites to a normal form. In particular, every
word rewrites to a normal form by starting from the identity normal form.
\end{theorem}

\todo{When \code{N.Completeness.lemma-coset-update} is finished, replace the
assumption by the completed theorem and state the final completeness corollary:
words modulo the symplectic relations are isomorphic to \(\Sp(2n,\Zp)\).}

\section{From Symplectic to Clifford-Mod-Scalar}
\label{sec:clifford}

The project next derives Clifford-mod-scalar relations from two simpler
presentations: a Pauli presentation and the symplectic presentation. The key
structure is a semidirect product. The Pauli generators \(X\) and \(Z\) form the
normal subgroup. The symplectic generators \(S\), \(H\), and \(CZ\) form the
quotient layer. Cross-relations describe the conjugation action of symplectic
generators on Pauli generators, for example how \(H\) swaps \(X\) and \(Z\), how
\(S\) maps \(X\) to \(XZ\), and how \(CZ\) maps an \(X\) on one wire to an
\(X\) together with a \(Z\) on the other.

This construction is formalized in \code{N.Clifford.SDProduct}. The combined
generator set is the disjoint union of Pauli generators and symplectic
generators. The combined relation is built using the generic semidirect-product
presentation constructor. The embedding of the symplectic presentation into the
Clifford-mod-scalar presentation is the right injection into this disjoint
union; the embedding of the Pauli presentation is the left injection. Thus a
symplectic relation can be reused as a Clifford relation by mapping every
symplectic generator through the embedding.

The comparison theorem \codepath{N.Clifford.Iso3.M.Theorem-SemiDirect-iso-Clifford}
proves that the semidirect-product presentation is isomorphic to the
Clifford-mod-scalar presentation. This theorem is the paper-level bridge
between the generic algebraic construction and the concrete Clifford relations.

\subsection{Simplifying the Clifford Presentation}

The first Clifford-mod-scalar presentation uses derived words involving
\(\mathsf{s}=S Z^{1/2}\) and multipliers. This is close to the semidirect-product
derivation but not ideal as a final rule set. The simplified presentation pushes
Pauli components to the right, uses basic \(S,H,CZ\) gates where possible, and
adds explicit \(Z\)-\(CZ\) commutation rules to avoid circular derivations.

The file \code{N.Clifford.Clifford-Mod-Scalars-Simplified} states the simplified
relations. The file \code{N.Clifford.Iso-Clifford-Simplified} proves that the
original and simplified presentations are isomorphic by the identity map. The
proof has two directions: every original relation is derivable in the simplified
presentation, and every simplified relation is derivable in the original
presentation.

A simple simplification example is \codepath{N.Clifford.Mg-Simplify.s-collect}:
for every \(k\), it collects the hidden Pauli part of \(\mathsf{s}^k\) and
rewrites it as
\[
  \mathsf{s}^k \approx S^k \cdot Z^{k/2}.
\]
This is the local shape of the entire simplification pass: expose the Pauli
components, move them to a canonical side, and compare the remaining
symplectic word. More complicated examples include the \code{final-semi-M*}
theorems in \code{N.Clifford.Mg-Simplify}, which replace multiplier relations
involving \(\mathsf{s}\) by collected \(S,H\)-words with rightmost Pauli tails.

\section{Adding Scalars as a General Group Extension}
\label{sec:scalars}

The semidirect-product construction above is split: quotient relations lift
without correction. Adding scalar phases requires a more general group-extension
presentation. Suppose \(N\) is a normal subgroup with presentation
\(\langle Y \mid S\rangle\), and the quotient \(G/N\) has presentation
\(\langle \overline{X}\mid\overline{R}\rangle\). Choosing representatives
\(X\subseteq G\) for the quotient generators gives two kinds of extra data.
First, each quotient relator may lift to a word in \(N\), giving a correction
term. Second, each quotient generator conjugates each kernel generator to a word
in \(N\).

The module \code{Presentation.Group-Extension2} formalizes this recipe. The
relations consist of kernel relations, corrected quotient relations, and
conjugation relations. The split semidirect product is recovered as the special
case where all quotient corrections are \(\epsilon\). Scalar phases are more
general: quotient equations that hold only up to a scalar become corrected
relations whose correction word lives in the scalar subgroup.

\todo{Instantiate this section with the scalar subgroup and the exact
correction words for the Clifford presentation once those files are finalized.}

\section{Other Case Studies}
\label{sec:examples}

The examples demonstrate that the library is not specialized to qudit Clifford
circuits.

\begin{table}[tbp]
\centering
\begin{tabular}{lll}
\toprule
File & Object & Framework feature \\
\midrule
\code{Examples/CliffordT1.agda} & one-qubit Clifford+\(T\) & amalgamation \\
\code{Examples/QutritCliffordT1.agda} & one-qutrit Clifford+\(T\) & amalgamation \\
\code{Examples/U33Di.agda} & \(U_3(\mathbb{Z}[1/2,i])\) & cosets/amalgamation \\
\bottomrule
\end{tabular}
\caption{Independent case studies using the same presentation machinery.}
\label{tab:examples}
\end{table}

Each case builds smaller presentations with normal forms, constructs coset
representatives or amalgamation data, and concludes with an isomorphism theorem
between the intended presentation and the presentation generated by the generic
construction.

The one-qubit Clifford+\(T\) example ends with
\codepath{Theorem-CliffordT1-iso-B*A**} in the final
\codepath{Examples.CliffordT1.CliffordT1} module. It builds a direct product of
cyclic presentations, extends it by coset normal forms, and then uses
\code{AmalDataNF}/\code{ANF} to obtain the amalgamated product normal form and
the final star-isomorphism theorem.

The one-qutrit Clifford+\(T\) example has two final isomorphism theorems named
\codepath{Theorem-CliffordT1-iso-B*A**}, one in
\codepath{Examples.QutritCliffordT1.CliffordT1} and one in the simplified
module. Both are built from chained coset normal forms and the same
amalgamation machinery.

The \(U_3(\mathbb{Z}[1/2,i])\) example ends with
\codepath{Examples.U33Di.Iso.Theorem-U33Di-iso-B*A**}. This case combines cyclic
products, \code{SnD}, coset normal forms, direct products, sugar products, and
amalgamation, making it the strongest evidence that the framework is not tied
to Clifford-specific syntax.

\section{Related Work}
\label{sec:related}

\paragraph{Circuit presentations and normal forms.}
Lafont's algebraic theory of Boolean circuits is an early systematic study of
generators and relations for circuit-like structures~\cite{Lafont2003}. In the
quantum setting, Selinger gave generators, relations, and normal forms for
\(n\)-qubit Clifford operators~\cite{Selinger2015}. Amy, Chen, and Ross gave a
finite presentation for CNOT-dihedral operators using normal forms and rewrite
systems~\cite{AmyChenRoss2018}. Presentations have also been studied for
larger exact-synthesis groups such as
\(U_4(\mathbb{Z}[1/\sqrt{2},i])\)~\cite{Greylyn2014} and for two-qubit
Clifford+\(T\) operators~\cite{BianSelinger2022}. These works are close in
proof shape to ours, but their proofs are carried out mathematically rather
than inside a proof assistant.

\paragraph{Qudit and Clifford completeness.}
The stabilizer and Clifford structures used here go back through work on
higher-dimensional fault-tolerant quantum computation and arbitrary-dimensional
Clifford operations~\cite{Gottesman1999,HostensDehaeneDeMoor2005}, as well as
structural characterizations of Clifford operators~\cite{Farinholt2014} and
normal-form decompositions of Clifford circuits~\cite{BravyiMaslov2021}. Recent
work gives complete and natural rule sets for multi-qutrit Clifford
circuits~\cite{LiMoscaRossVanDeWeteringZhao2025Qutrit}; the present
formalization generalizes this style of normal-form reasoning to all odd prime
dimensions at the symplectic level and mechanizes the algebra in \Agda. The
companion multi-qudit Clifford development for all odd primes uses the same
normal boxes and relation-simplification
strategy~\cite{BianLiRossVandeWeteringZhaoQudit}.

\paragraph{Beyond Clifford.}
Matsumoto and Amano introduced normal forms for single-qubit Clifford+\(T\)
circuits~\cite{MatsumotoAmano2008}; Giles and Selinger clarified this normal
form and its exact-synthesis consequences~\cite{GilesSelinger2013}. Glaudell,
Ross, and Taylor gave canonical forms for single-qutrit Clifford+\(T\)
operators~\cite{GlaudellRossTaylor2018}, and Jain, Kalra, and Prakash extended
normal-form ideas to single-qudit Clifford+\(T\)
operators~\cite{JainKalraPrakash2020}. Complete equational theories for
broader circuit fragments include the general quantum-circuit theory of Clement
et al.~\cite{ClementEtAl2022} and the real-Clifford+\(CH\) theory of
Clement~\cite{Clement2026RealCH}. These results motivate our future-work goal
of formalizing the first \(n\)-qubit universal completeness theorem.

\paragraph{Graphical calculi.}
ZX-calculi provide complete graphical languages for several fragments of
quantum mechanics, including stabilizer fragments and universal
fragments~\cite{CoeckeDuncan2011,Backens2014,JeandelPerdrixVilmart2018}. The
ZX tradition is complementary to our circuit-presentation approach: ZX systems
reason in a richer graphical language, while this work focuses on syntactic
presentations of circuit generators and relations. Verified ZX efforts such as
VyZX show a related proof-assistant direction for graphical languages.

\paragraph{Proof assistants for quantum circuits.}
QWIRE embeds quantum circuits in Coq as a core circuit
language~\cite{PaykinRandZdancewic2017QWIRE} and supports formal verification
of circuit equivalences~\cite{RandPaykinZdancewic2018QWIRE}; ReQWIRE extends
this style of reasoning to reversible quantum circuits~\cite{RandPaykinLeeZdancewic2019ReQWIRE}.
SQIR and VOQC develop a Coq-based intermediate representation and verified
optimizer~\cite{HietalaEtAl2021VOQC}; and VyZX formalizes ZX-calculus reasoning
in Coq~\cite{LehmannCaldwellRand2022VyZX,LehmannEtAl2023VyZX}. Our work
differs in its focus on presentations, normal forms, and completeness of
circuit relations rather than semantic preservation of compiler optimizations.

\paragraph{Circuit languages and interoperability.}
Quipper, OpenQASM, and Qiskit provide practical circuit description and
compilation ecosystems~\cite{Selinger2004QPL,GreenEtAl2013Quipper,CrossEtAl2022OpenQASM3,Wille2019Qiskit}.
Modern high-level languages such as Silq address programming abstractions like
safe uncomputation~\cite{BichselEtAl2020Silq}.
Interfacing the \Agda framework with such languages would allow external circuit
programs to be imported into a certified presentation, normalized, and exported
with proof-producing rewrites.

\section{Future Work}

First, the most direct extension is a formalization of normal forms and complete
relations for Galois qudit Clifford circuits. Second, the current examples use
at most binary gates; extending the framework to ternary and general \(n\)-ary
gates would make it more suitable for universal circuit families and graphical
calculi with arbitrary arities. Third, practical impact requires interfaces with
circuit description languages such as Quipper, Qiskit, and OpenQASM. Finally,
we plan to formalize the complete equational theory for real-Clifford+\(CH\)
circuits, the first completeness result for an \(n\)-qubit universal circuit
fragment without parameterized gates or ancillas.

\section{Artifact Scope}
\label{sec:artifact}

The repository contains both core developments and exploratory modules. For a
POPL artifact, the trusted slice should be made explicit. The core slice should
include the word/presentation framework, the symplectic normal-form modules, the
box-relation derivations used by completeness, the Clifford-mod-scalar and
simplified Clifford modules, the group-extension machinery, and the three
example developments. The final artifact should either complete or exclude the
\code{Test} modules in \code{Examples/CliffordT1.agda} and
\code{Examples/QutritCliffordT1.agda}, which currently contain active holes;
\code{Examples/U33Di.agda} appears clean. Exploratory files with holes,
archived files, and older dimension-specific prototypes should be excluded from
the artifact claim unless they are completed and added to the build.

\bibliographystyle{ACM-Reference-Format}
\bibliography{popl2027}

\end{document}
